\section{Minimizing Total Wait Time}\label{sec: min wait time}
\doHeader{Minimizing Total Wait Time}

\paragraph{Problem Definition}
\begin{mylist}
\item[\bf input:] a sequence $T=(t_1, t_2, \ldots, t_n)$ of $n$ \emph{tasks},
  each with an associated \emph{length} $\ell(t_i)$.
\item[\bf output:] a permutation $T'=(t'_1, t'_2, \ldots, t'_n)$ of $T$ minimizing the \emph{total wait time},
  denoted  $W(T')$ and defined as $W(T') = \sum_{i=1}^n \sum_{j=1}^i \ell(t'_j)$.
\end{mylist}
Intuitively, we must choose an order in which to do the tasks
so as to minimize the sum of their respective times until completion,
assuming the $j$th task we do (task $t'_j$)
starts when the $(j-1)$st task finishes,
and then takes time $\ell(t'_j)$ equal to the task's length.

We start with some observations that will be useful for developing an algorithm.
\begin{observation}\label{obs: wait time 1} 
  % Fix any instance $T=(t_1, \ldots, t_n)$ and permutation $T'=(t'_1, \ldots, t'_n)$.
  The wait time $W(T')$ can also be expressed as
  \[
    W(T') = \sum_{j=1}^n (n-j+1)\, \ell(t'_j).
    % = n\, \ell(t'_1) + (n-1) \ell(t'_2) + \cdots + 2\, \ell(t'_{n-1}) + \ell(t'_n).
  \]
\end{observation}
The observation holds by considering the sum that defines $W(T')$:
the $j$th task $t'_j$ contributes its length $\ell(t'_j)$
to its own completion time, and to the completion times of the $n-j$ following jobs.

Next we consider how swapping two adjacent tasks changes the wait time.
\begin{observation}\label{obs: wait time 2} 
  For any permutation $T'=(t'_1, \ldots, t'_n)$,
  and any $j$ with $1\le j \le n-1$,
  let $T''=(t''_1, \ldots, t''_n)$ be the permutation obtained from $T'$ by swapping
  the $j$th job $t'_j$ with the following job $t'_{j+1}$.
  That is, $t''_j = t'_{j+1}$, $t''_{j+1} = t'_j$,
  and $t''_i = t'_i$ for all $i\not\in \{j, j+1\}$.

  Then $W(T'') = W(T') + \ell(t'_{j+1}) - \ell(t'_j)$.
\end{observation}
The observation follows from the previous one.
(Swapping increases the coefficient of $\ell(t'_{j+1})$ by 1,
and decreases the coefficient of  $\ell(t'_{j})$ by 1.)

We'll use the observation to show that $t'_1$ should be the shortest task.
The next observation will then allow us to decouple the problem of choosing the first task $t'_1$
from the problem of ordering the remaining $n-1$ tasks, $(t'_2, \ldots, t'_n)$.
\begin{observation}\label{obs: wait time 3} 
  For any permutation $T'=(t'_1, \ldots, t'_n)$,
  let
  $R'=(r'_1, r'_2, \ldots, r'_{n-1}) = (t'_2, t'_3, \ldots, t'_{n})$
  be the permutation that $T'$ induces on the $n-1$ tasks other than $t'_1$.
  Then
  \[
    W(T') = n\, \ell(t'_1) + W(R').
  \]
\end{observation}
Note that $W(R')$ is the total wait time of the permutation $R'$ of the remaining $n-1$ tasks,
considered as an independent instance with just those $n-1$ tasks.
(Observation~\ref{obs: wait time 3} can be verified by using Observation~\ref{obs: wait time 1}
to expand $W(T')$ and $W(R')$, then using $r'_j = t'_{j+1}$.)

\paragraph{Algorithm.}
Now we consider the first greedy choice.
We use Observation~\ref{obs: wait time 2} to show that
we can put any shortest task first:

\begin{lemma}\label{lemma: wait time 1}
  Consider any instance $T=(t_1, \ldots, t_n)$ with $n\ge 1$.
  Let $t_m$ be any task of minimum length, so $\ell(t_m) = \min_i \ell(t_i)$.
  Then there is an optimal solution $T'$ that starts with job $t_m$.
  % That is, $T'$ is of the form $T'=(t_m, r'_1, r'_2, \ldots, r'_{n-1})$,
  % where  $R'=(r'_1, r'_2, \ldots, r'_{n-1})$ is some permutation of the remaining $n-1$ tasks.
\end{lemma}
\begin{longFormProof}[Proof of Lemma~\ref{lemma: wait time 1}.]
  \step Consider any instance $T$ with $n\ge 1$ and let $t_m$ be any task of minimum length.
  \step Let $T'$ be an optimal permutation (there is one, as the number of permutations is finite).
  \begin{block}
    \step \underline{Case 1.} Suppose $T'$ starts with task $t_m$ (that is, $t'_1 = t_m$).
    \step Then there is an optimal solution (namely $T'$) that starts with task $t_m$.
  \end{block}
  \begin{block}
    \step \underline{Case 2.} Otherwise task $t_m$ is not first in $T'$.
    \step Consider bringing task $t_m$ to the front by repeatedly swapping it with the task before it.
    \step By Observation~\ref{obs: wait time 2} (and $\ell(t_m) \le \ell(t_j)$ for all $j$)
    each swap does not increase the total wait time.
    \step So the resulting permutation (which starts with task $t_m$) is also optimal.
  \end{block}
\end{longFormProof}

By Lemma~\ref{lemma: wait time 1}, it is safe to choose any minimum-length task $t_m$ to go first.
Next we use Observation~\ref{obs: wait time 3},
to show that, once the first task is chosen as $t_m$,
we can simply recurse on the remaining $n-1$ tasks.
That is, any optimal permutation of those tasks (considered as an independent instance with $n-1$ tasks)
will work.

\begin{lemma}\label{lemma: wait time recurse}
  Consider any instance $T=(t_1, t_2, \ldots, t_n)$ with $n\ge 1$.
  Let $t_m$ be any task of minimum length, so $\ell(t_m) = \min_i \ell(t_i)$.
  Let $R=(r_1, r_2, \ldots, r_{n-1})$ be the $n-1$ remaining tasks (other than $t_m$).
  Let $R'=(r'_1, r'_2, \ldots, r'_{n-1})$ be any optimal permutation of those $n-1$ tasks.

  Then the permutation $T'=(t_m, r'_1, r'_2, \ldots, r'_{n-1})$ is optimal for $T$.
\end{lemma}
\begin{longFormProof}
  \step Consider any such $T$, $m$, $R$, $R'$, and $T'$.

  \step Let $T^*=(t^*_1, t^*_2, \ldots, t^*_{n})$ be an optimal permutation of $T$ with $t^*_1=t_m$
  ($T^*$ exists by Lemma~\ref{lemma: wait time 1}).

  \step Let $R^*=(t^*_2, \ldots, t^*_n)$ be the permutation that $T^*$ induces on the tasks other than $t_m$.

  \step Since $R'$ is an optimal permutation of those $n-1$ tasks, $W(R') \le W(R^*)$.

  \step Using  Observation~\ref{obs: wait time 3}, then $W(R') \le W(R^*)$, then Observation~\ref{obs: wait time 3} again,
  \[
    W(T') \,=\, n\, \ell(t_m) + W(R') \,\le\, n\, \ell(t_m) + W(R^*) \,=\, W(T^*).
  \]

\end{longFormProof}

This gives us the following algorithm:

\begin{myframed}
  \begin{steps}
  \item[{\textsf{rMinWaitTime}$(t_1, t_2, \ldots, t_n)$:}]
    \step if $n = 0$: return the empty sequence
    \step let $t_m$ be a job of minimum length
    \step let $(r_1, r_2, \ldots, r_{n-1})$ be the $n-1$ jobs other than $t_m$, in any order
    \step let $(r'_1, r'_2, \ldots, r'_{n-1}) = \textsf{rMinWaitTime}(r_1, r_2, \ldots, r_{n-1})$
    \step return $T'=(t_m, r'_1, \ldots, r'_{n-1})$
  \end{steps}
\end{myframed}

\paragraph{Correctness}
\begin{lemma} The algorithm is correct.
\end{lemma}
\begin{shortFormProof}
  The proof is by induction on $n$.
  For the base case $n=0$, the algorithm is clearly correct.
  Consider any $n\ge 1$.
  Assume inductively that the algorithm is correct on all inputs of length $n-1$.
  Consider executing the algorithm on any input $T=(t_1, t_2, \ldots, t_n)$ of length $n$.
  Let $t_m$, $(r_1, r_2, \ldots, r_{n-1})$, $(r'_1, r'_2, \ldots, r'_{n-1})$, 
  and $T'$ be as computed in the algorithm.
  (See the algorithm for their definitions.)
  By the inductive assumption $(r'_1, r'_2, \ldots, r'_{n-1})$ is optimal for $(r_1, r_2, \ldots, r_{n-1})$.
  Hence, by Lemma~\ref{lemma: wait time recurse}, $T'$ is optimal for $T$.
\end{shortFormProof}

\paragraph{Run time.}
Clearly the algorithm can be implemented to run in $O(n^2)$ time.
But we can do better.  Observe that the algorithm simply sorts the jobs
by increasing length.  This can be done via Mergesort in $O(n\log n)$ time.
So  the algorithm can be implemented to run in $O(n\log n)$ time.

The next theorem summarizes what we've figured out:
\begin{theorem}
  There is an $O(n\log n)$-time algorithm for \prob{Minimizing Total Wait Time}.
\end{theorem}

Exercise~\ref{exercise: min wait time} % (Section~\ref{sec: exercises})
asks you to extend this result to a more general
% [review 2026-09-27] fix: weights must be positive (not just non-negative), since the exercise orders jobs by $\ell_i/w_i$, which is undefined for $w_i = 0$
problem where each task also has a positive \emph{weight},
and the goal is to minimize the \emph{weighted} sum of completion times.

\subsection{Exercises}

\paragraph{Exercise without solution}

\exercise\label{exercise: min wait time}
Consider the following generalization of the problem in Section~\ref{sec: min wait time}.
\begin{mylist}
\item[\bf input:] an integer $n\ge 0$, and for each $i\in\{1,\ldots,n\}$,
  % [review 2026-09-27] fix: weight must be positive so that $\ell_i/w_i$ (used below) is defined
  a \emph{job length} $\ell_i$ and positive \emph{weight} $w_i$.
\item[\bf output:] a permutation $\pi$ of the $n$ jobs minimizing the \emph{weighted sum of completion times},
  % [review 2026-09-27] fix: objective attached weight $w_i$ to position $i$ instead of to the job $\pi(i)$ in that position (as written, ordering by $\ell_i/w_i$ is not optimal); changed to $w_{\pi(i)}$
  defined as $\sum_{i=1}^n w_{\pi(i)} \sum_{j=1}^i \ell_{\pi(j)}$.
\end{mylist}
Prove carefully that ordering the jobs by increasing $\ell_i/w_i$ works
(generalizing the algorithm in Section~\ref{sec: min wait time}).
% [review 2026-09-27] fix: "even when $(\ell_i,w_i) = (\ell_j,w_j)$" should be $\ne$ (ties in the ratio occur with unequal pairs, e.g. (1,2) and (2,4))
Make sure your proof correctly handles instances with ties, where $\ell_i/w_i = \ell_j/w_j$ for some jobs $i$ and $j$.  This can happen even when $(\ell_i, w_i) \ne (\ell_j, w_j)$!


\subsection{External sources on \prob{Minimizing Total Wait Time}}

\begin{itemize}
\item 
  \href{http://jeffe.cs.illinois.edu/teaching/algorithms/book/04-greedy.pdf}{JE Chapter 4} Section 4.1
  (As of August 2019,
  Erickson's presentation has a subtle error.
  In our terminology, his Lemma 4.1 states that
  \emph{any permutation that orders the tasks by increasing length is optimal}.
  But what his proof actually shows is the converse:
  \emph{any permutation that doesn't order the tasks by increasing length isn't optimal}.
  These are not the same!  We've filed a bug report.)
  
\end{itemize}

% \section{Minimizing Total Wait Time}\label{sec: min wait time}
% \doHeader{Minimizing Total Wait Time}

% \paragraph{Problem Definition.}
% \begin{mylist}
% \item[\bf input:] a sequence $\ell=(\ell_1, \ell_2, \ldots, \ell_n)$ of non-negative numbers, called \emph{job lengths}
% \item[\bf output:] a permutation $\pi$ of the $n$ jobs minimizing the \emph{total wait time},
%   defined as $\sum_{i=1}^n \sum_{j=1}^i \ell_{\pi(j)}$.
% \end{mylist}
% Intuitively, there are $n$ jobs $1, 2, \ldots, n$.
% We must choose a permutation $\pi(1), \pi(2), \ldots, \pi(n)$ of those jobs
% so as to minimize the sum of their respective wait times (time until completion),
% assuming the $j$th job we do (job $\pi(j)$)
% starts when the $(j-1)$st job finishes,
% and takes additional time $\ell_{\pi(j)}$.

% \paragraph{Algorithm.}
% Let's try the ``make one choice then recurse'' approach.
% For this problem, a natural first choice is to commit to placing the shortest job first.
% The algorithm does this, then recurses on the remaining jobs:

% \begin{myframed}
%   \begin{steps}
%   \item[{\textsf{rMinWaitTime}$(\ell_1, \ell_2, \ldots, \ell_n)$:}]
%     \step if $n = 0$: return the empty sequence
%     \step let $\ell_j$ be a job of minimum length
%     \step let $\ell'_1, \ell'_2, \ldots, \ell'_{n-1}$ be the jobs other than $\ell_j$
%     \step return $\ell_j$ followed by \textsf{rMinWaitTime}$(\ell'_1, \ell'_2, \ldots, \ell'_{n-1})$
%   \end{steps}
% \end{myframed}

% This can be implemented in $O(n\log n)$ time using, say, Mergesort.

% Note that if some jobs have equal lengths, the output depends on how ties are broken.

% We can also think of this algorithm as follows:
% choose the shortest job to go first,
% then recurse on the remaining jobs.
% ``Choose one piece then recurse''
% is a common template for many greedy algorithms.

% \paragraph{Correctness.}
% We first calculate how swapping two adjacent changes the wait time.
% \begin{observation}\label{obs: wait time} 
%   Fix any input $\ell_1, \ell_2, \ldots, \ell_n$ and permutation $\pi$.
%   Let $\pi'$ be the permutation obtained from $\pi$ by swapping
%   two jobs $\pi(j)$ and $\pi(j+1)$, for some given $j$.
%   That is, $\pi'(j) = \pi(j+1)$, $\pi'(j+1) = \pi(j)$,
%   and $\pi'(i) = \pi(i)$ for all $i\ne j$.

%   Then the total wait time for $\pi'$ equals the total wait time for $\pi$ plus
%   $\ell_{\pi(j+1)} - \ell_{\pi(j)}$.
% \end{observation}
% The observation holds by calculation.
% (After swapping,
% job $\pi(j)$ now has to wait for job $\pi(j+1)$,
% increasing the total wait time by $\ell_{\pi(j+1)}$.
% But job $\pi(j+1)$ no longer has to wait for job $\pi(j)$,
% decreasing the total wait time by $\ell_{\pi(j)}$.)

% Now we show that the algorithm is correct:

% \begin{lemma}\label{lemma: ordered is opt}
%   For any input $\ell_1, \ell_2, \ldots, \ell_n$, any permutation $\pi$
%   that orders the jobs by increasing length is optimal.
%   That is, the algorithm is correct.
% \end{lemma}
% \begin{longFormProof}
%   \step Consider any input $\ell$ and permutation $\pi$ that orders the jobs by increasing length.
%   \step Let $\pi^*$ be an optimal permutation (there is at least one, as the number of permutations is finite).
%   \step Then $\pi^*$ also orders the jobs by increasing length.
%   (Otherwise it has a job followed by a shorter job.
%   By Observation~\ref{obs: wait time} swapping them
%   decreases the wait time, contradicting that $\pi^*$ is optimal).
%   \step So $\pi$ and $\pi^*$ both order the jobs by increasing length.
%   \step So $\pi$ can be obtained from $\pi^*$ by repeatedly swapping
%   adjacent, equal-length jobs.
%   \step By Observation~\ref{obs: wait time}, each such swap preserves the wait time.
%   \step So the two permutations give the same total wait time, and $\pi$ is also optimal.
% \end{longFormProof}

% We have shown the following result:
% \begin{theorem}
%   There is an $O(n\log n)$-time algorithm for Minimizing the Total Wait Time.
% \end{theorem}

% Exercise~\ref{exercise: min wait time} (Section~\ref{sec: exercises})
% asks you to extend this result to a more general
% problem where each job also has a non-negative \emph{weight},
% and the goal is to minimize the \emph{weighted} sum of completion times.

% \paragraph{Other resources on \prob{Minimizing Total Wait Time}}

% \begin{itemize}
% \item 
%   \href{http://jeffe.cs.illinois.edu/teaching/algorithms/book/04-greedy.pdf}{JE Chapter 4} Section 4.1
%   (As of August 2019,
%   Erickson's presentation subtly glosses over the issue of equal-length jobs.
%   In our terminology, his Lemma 4.1 states that
%   \emph{any permutation that orders the jobs by increasing length is optimal}
%   (equivalent to the second lemma above).
%   But what his proof actually shows is the converse:
%   \emph{any optimal permutation orders the jobs by increasing length},
%   (equivalent to the first lemma above).
%   We've filed a bug report!)
  
% \end{itemize}

%%% Local Variables:
%%% mode: latex
%%% TeX-master: "greedy_algorithms"
%%% End:
